\section{The Temporal Attention Neuron}\label{sec:model}

Traditional spiking neuron models such as the leaky integrate-and-fire (LIF) neuron are built on continuous-time ordinary differential equations with strictly Markovian dynamics. In contrast, the \textbf{Temporal Attention Neuron (TAN)} introduced here is a discrete-time dynamical system that maintains a short local history (a window of length $W$) of past inputs, yet still relies only on local update rules (leakage and threshold) at each step. The model combines a Boltzmann-style attention kernel with neuromodulatory gain control inside a single neuron.

\subsection{Model Formulation}

The core membrane dynamics of a TAN follow a leaky integrator, where a synaptic current $A_t$ drives the membrane potential $h_t \in \mathbb{R}$:
\begin{equation}\label{eq:membrane}
h_t = \lambda \, h_{t-1} + A_t, \qquad
y_t = \Theta(h_t - \theta),
\end{equation}
where $\lambda \in (0,1)$ is the leak factor (determined by the membrane RC time constant and guaranteeing bounded, stable decay), $\Theta(\cdot)$ is the Heaviside step function, $\theta$ is the firing threshold, and $y_t \in \{0,1\}$ is the output spike. After a spike ($y_t = 1$), the membrane may optionally be reset to a resting potential $h_{\text{rest}}$ for biophysical consistency.

This discrete-time formulation is analogous to the standard LIF neuron and rests on classical foundations of neural modeling~\cite{hodgkin1952,abbott1999}.

\subsection{Deviation Signal and Asymmetric Rectification}

At each time step $t$, the TAN maintains a finite \textbf{temporal receptive field} consisting of the last $W$ inputs:
\begin{equation}
\window = [\,x_{t-W+1},\; x_{t-W+2},\; \dots,\; x_t\,],
\end{equation}
and computes the local mean $\mu_t = \frac{1}{W}\sum_{i=1}^W \mathcal{X}_{t,i}$. Note that this mean includes the current observation $x_t$. The \textbf{deviation signal} $S_t$ is then defined as
\begin{equation}
S_t = x_t - \mu_t.
\end{equation}

Thus $S_t$ measures how much the current input deviates from the local temporal-window mean. We refer to this as a \textit{prediction-error-inspired signal} or \textit{local surprise}, by analogy with predictive coding theory~\cite{friston2010,rao1999,srinivasan1982,friston2017}. However, we note that $S_t$ is not a prediction error generated by a predictor based exclusively on past observations; it is the deviation of the current observation from a window mean that includes $x_t$ itself. A full predictive-coding architecture would require a separate generative model and a strictly past-only predictor.

To allow the neuron to ignore small fluctuations below the environmental noise floor, we further apply asymmetric rectification with a noise-tolerance dead-zone $\tau \geq 0$:
\begin{equation}
S_t' = \max(0,\, S_t - \tau).
\end{equation}

This nonlinearity provides a computational analogue of asymmetric response or gating. In biological systems, chemical transmitter concentrations have no ``negative absolute value'', and the opening of ion channels is direction-dependent~\cite{kandel2000}. In the macroscopic Transformer architecture, every self-attention layer is followed by a feed-forward network (FFN) equipped with a nonlinear activation function such as ReLU or GELU~\cite{vaswani2017}; the asymmetric rectification in TAN can be viewed as a single-neuron analogue of this FFN nonlinearity, which is needed to transform bidirectional deviation signals into unidirectional adaptive responses (see~\S\ref{sec:evaluation}).

\subsection{Temporal Attention}

The TAN employs an attention-like kernel. First, the rectified deviation signal and the history window are projected into query, key, and value vectors:
\begin{equation}
Q_t = \phi_q(S_t'), \qquad K_{t,i} = \phi_k(\mathcal{X}_{t,i}), \qquad V_{t,i} = \phi_v(\mathcal{X}_{t,i}),
\end{equation}
where $\phi_q, \phi_k, \phi_v$ are (possibly learnable) linear or nonlinear synaptic mappings. The query $Q_t$ is driven by the rectified deviation signal $S_t'$ rather than by the raw input $x_t$. When the surprise signal falls below the gating threshold, $S_t' = 0$ and thus $Q_t = 0$. In this case, all products $Q_t \cdot K_{t,i} = 0$, so the attention weights become uniform: $\alpha_{t,i} = 1/W$ for all $i$. The attention distribution does not vanish; rather, it becomes uniform. The effective suppression of the neuron's response comes from the subsequent gated amplitude $A_t = g(S_t') \cdot C_t$, because $g(0) = 0$ ensures that no effective current is injected into the membrane even when a uniform context is retrieved.

Next, a \textbf{compatibility score} is computed between the current query and each past key:
\begin{equation}
E_{t,i} = Q_t \cdot K_{t,i}.
\end{equation}

The attention weights are then given by the normalized Boltzmann (Gibbs) distribution
\begin{equation}
\alpha_{t,i} = \frac{\exp(\beta \, E_{t,i})}{\sum_{j=1}^W \exp(\beta \, E_{t,j})},
\end{equation}
where $\beta > 0$ is an inverse-temperature parameter that controls the sharpness of the attention. Since $\exp(+\beta E_{t,i})$ is monotonically increasing in $E_{t,i}$, larger $\beta$ focuses weight on the \textit{maximum-compatibility} (highest $E_{t,i}$) historical inputs. The context vector is the attention-weighted sum of the values:
\begin{equation}
C_t = \sum_{i=1}^W \alpha_{t,i} \, V_{t,i}.
\end{equation}

The weights $\alpha_{t,i}$ take the form of a Gibbs distribution over the compatibility scores $E_{t,i}$~\cite{jaynes1957,boltzmann1868,gibbs1902}. We emphasise that $E_{t,i}$ should be understood as a \textit{compatibility score} (higher is better), not as a physical energy to be minimised. The neuron allocates its synaptic influence toward high-compatibility (best-matching) historical inputs. This provides a connection to statistical-mechanical interpretations of deep learning~\cite{bahri2020}, although we do not claim that TAN is mathematically equivalent to a Transformer.

\subsection{Gating Mechanism}

The raw context $C_t$ is not injected directly into the membrane. Instead, it is modulated by a nonlinear gain function of the rectified deviation signal:
\begin{equation}
A_t = g(S_t') \; C_t,
\end{equation}
where $g: \mathbb{R}^+ \to [0,1]$ is a monotonically increasing function satisfying $g(0) = 0$ and saturating to $1$ for large arguments. Two canonical choices are the sigmoid $g(x) = (1 + e^{-k(x-x_0)})^{-1}$ and the hyperbolic tangent $g(x) = \tanh(x)$, which is used throughout this paper for its origin-preserving property ($g(0) = 0$ exactly, with no offset).

This multiplicative gating is computationally analogous to the effect of global neuromodulatory systems such as the locus coeruleus--norepinephrine (LC-NE) pathway, which is known to implement adaptive gain control~\cite{aston2005,friston2010}. Only when the deviation magnitude $|S_t'|$ is sufficiently large does the gate open and allow the context $C_t$ to contribute to the synaptic current $A_t$. When the stimulus becomes predictable ($x_t \approx \mu_t$ and $S_t' \approx 0$), the gate shuts ($g \approx 0$) and the current vanishes, effectively decoupling the neuron from the past context. The gate thus implements a ``switch'' between attentive processing and quiet silence at the single-cell level.

\subsection{Membrane Dynamics and Complete System}

Assembling all components, the TAN dynamics are described by the following closed set of equations:
\begin{align}
\window &= [x_{t-W+1},\dots,x_t], \quad \mu_t = \tfrac{1}{W}\textstyle\sum_{i=1}^W \mathcal{X}_{t,i}, \quad S_t = x_t - \mu_t, \notag \\
S_t' &= \max(0,\, S_t - \tau), \notag \\
Q_t &= \phi_q(S_t'), \quad
K_{t,i} = \phi_k(\window_i), \quad
V_{t,i} = \phi_v(\window_i), \notag \\
\alpha_{t,i} &= \frac{\exp(\beta \, Q_t \cdot K_{t,i})}{\sum_{j=1}^W \exp(\beta \, Q_t \cdot K_{t,j})}, \quad
C_t = \sum_{i=1}^W \alpha_{t,i} V_{t,i}, \notag \\
A_t &= g(S_t') \, C_t, \qquad
h_t = \lambda \, h_{t-1} + A_t, \qquad
y_t = \Theta(h_t - \theta). \notag
\end{align}

The system is entirely specified by local, time-step updates and nonlinear transformations, without any external recurrence or learning in its minimal form. In our default configuration, the synaptic mappings are linear scalars $\phi_q(x) = W_q x$, $\phi_k(x) = W_k x$, $\phi_v(x) = W_v x$ with $W_q = 2.0$, $W_k = W_v = 1.0$, the inverse temperature $\beta = 1.0$, the leak $\lambda = 0.5$, and the threshold $\theta = 0.5$.

\subsection{Theoretical Properties}

The TAN model possesses three theoretically interesting characteristics that distinguish it from both classical LIF neurons and standard Transformer blocks.

\textbf{Bridging Markovian and non-Markovian dynamics.} The membrane update $h_t = \lambda h_{t-1} + A_t$ retains a strictly Markovian leaky integration, which guarantees boundedness and stability. However, the current $A_t$ is computed via an attention kernel that spans a non-Markovian window $\window$. TAN is thus a hybrid dynamical system: it superimposes a long-range temporal receptive field onto a local, Markovian substrate. This allows a single neuron to respond to temporal patterns extending over multiple time steps without requiring any auxiliary recurrent connections.

\textbf{Attracting fixed point under constant input.} In many classical spiking neuron models, habituation requires additional adaptation currents or slow negative feedback~\cite{benda2003,neftci2019}. In TAN, no extra ``habituation module'' is needed. When the input becomes constant ($x_t = c$), the deviation $S_t = x_t - \mu_t$ decays to zero. Consequently the gate $g(S_t') \to 0$ and the synaptic current $A_t = 0$. The membrane potential then simply follows $h_t = \lambda h_{t-1}$, decaying until it falls below threshold and spiking ceases. The zero-membrane state is an attracting fixed point under sustained constant input, providing a dynamical mechanism for habituation. This is formally proven in~\S\ref{sec:dynamics}.

\textbf{Attention as a compatibility-weighted retrieval.} The attention weights $\alpha_{t,i}$ take the form of a Gibbs distribution over the compatibility scores $E_{t,i}$~\cite{jaynes1957,boltzmann1868,gibbs1902}. The neuron's selection of a context vector $C_t$ can be viewed as a compatibility-weighted retrieval from recent history. This provides a connection to statistical-mechanical interpretations of deep learning~\cite{bahri2020}, although we do not claim formal mathematical equivalence between TAN and a Transformer.

% ============================================================
% 4. Emergent Dynamics of a Single TAN Neuron
% ============================================================
